\section{Core Primitives}
\label{sec:primitives}

This section names the artifacts that appear in the rest of the paper and fixes notation.

\subsection{Peers and the Ring}

A \emph{peer} is an instance of the Freenet Core running on some host. Each peer $p$ has a \emph{location} $\ell(p) \in [0, 1)$. Locations live on a one-dimensional ring with wrap-around distance
\[
  d(x, y) \;=\; \min\!\left\{|x - y|,\ 1 - |x - y|\right\}.
\]
A peer's location is derived deterministically from its external network address; the exact derivation and its security implications are discussed in \cref{sec:trust}. Each peer maintains a set of \emph{neighbor} connections, bounded between configurable minimum and maximum sizes (typically 25 and 200). Neighbor relations are reorganized over time by the topology rule described in \cref{sec:routing}.

\subsection{Contracts}

A \emph{contract} $\contract$ is a pair $(\mathrm{code}, \mathrm{params})$, where $\mathrm{code}$ is a WebAssembly module conforming to the platform's contract interface and $\mathrm{params}$ is an opaque byte string fixed at instantiation. The contract's \emph{key} is
\[
  k(\contract) = H\bigl(H(\mathrm{code}) \,\|\, \mathrm{params}\bigr)
\]
for a cryptographic hash $H$. The two-stage form lets a client that already holds $H(\mathrm{code})$ derive the key for a new contract from new parameters without re-hashing the full WebAssembly module, which may be several megabytes. The contract's \emph{location} $\ell(\contract) \in [0, 1)$ is derived from $k(\contract)$ in the same way as a peer location. Each contract has a \emph{state} $\sigma \in \statespace_\contract$, replicated across peers near $\ell(\contract)$ on the ring. The contract's code defines $\statespace_\contract$, the validity predicate $\valid_\contract$, the merge operation $\merge_\contract$, and the summary/delta synchronization functions (all detailed in \cref{sec:updates}). Contracts are public: their state is readable by any peer that can route to $\ell(\contract)$.

\subsection{Delegates}

A \emph{delegate} is a WebAssembly module that runs on a single user's device, inside the local Freenet Core. It holds private \emph{secret state} (sandboxed by the core) and responds to messages from user interfaces, contracts, and other delegates. The core attaches a verified sender identifier to every inbound message, so a delegate can apply policy based on sender identity in addition to message content. Delegates are local; they do not appear on the ring. The trust boundary they enforce and the cross-device replication story are described in \cref{sec:trust}.

\subsection{User Interfaces}

A \emph{user interface} is a web frontend (HTML, JavaScript, WebAssembly) delivered by a contract whose state encodes the application bundle. User interfaces run in the browser and communicate with the local core over a WebSocket API. They reach contracts and delegates only through that local core. They are not a distinct kind of network-visible artifact: from the platform's perspective, the bundle is just contract state with a particular MIME type.

% Per-peer architecture: what runs in and around one peer's process.
% Horizontal layout: browser (UI bundle) on left -- WebSocket -- Freenet
% Core (contracts + delegates) in middle -- UDP -- rest of network on right.
% Transport sits as a sub-box at the bottom of the Core.
\begin{figure}[t]
\centering
\begin{tikzpicture}[
    wasm/.style={draw=black!60, rounded corners=2pt, fill=orange!18,
                 minimum width=1.9cm, minimum height=0.65cm,
                 font=\small, align=center},
    ui/.style={draw=black!60, rounded corners=2pt, fill=purple!12,
               minimum width=2.4cm, minimum height=1.0cm,
               font=\small, align=center},
    transport/.style={draw=black!60, rounded corners=2pt, fill=green!12,
                      minimum width=5.0cm, minimum height=0.5cm,
                      font=\scriptsize\itshape, align=center},
    core/.style={draw=black!70, rounded corners=4pt, thick,
                 minimum width=5.3cm, minimum height=2.6cm,
                 fill=blue!4},
    netcloud/.style={draw=black!50, dashed, cloud, cloud puffs=12,
                     cloud puff arc=120, minimum width=2.5cm,
                     minimum height=1.3cm, font=\scriptsize\itshape,
                     fill=gray!8},
    arr/.style={->, >=Stealth, thick, black!70},
    label/.style={font=\scriptsize, gray!40!black},
    tag/.style={font=\scriptsize\itshape, gray!40!black},
    >=Stealth,
]
  % --- Browser on the left ---
  \node[ui] (ui) at (0,0) {UI bundle\\\scriptsize(HTML+JS+WASM)};
  \node[label, anchor=south] at (ui.north) {browser};

  % --- Freenet Core in the middle ---
  \node[core] (core) at (5.4,0) {};
  \node[font=\small\bfseries, anchor=north]
       at (core.north) [yshift=-3pt] {Freenet Core};
  \node[label, anchor=north] at (core.north) [yshift=-19pt] {(Rust, async)};

  % Contracts and delegates as a row inside the core.
  \node[wasm] (c1) at (4.4,-0.1)  {contract};
  \node[wasm] (d1) at (6.4,-0.1) {delegate};

  % Transport row attached to the bottom of the core.
  \node[transport, anchor=south]
       at ([yshift=4pt] core.south) (t)
       {transport: UDP, NAT hole-punch, X25519+AEAD};

  % --- Rest of network on the right ---
  \node[netcloud, right=14mm of core.east] (net) {rest of network};

  % --- Wires ---
  \draw[arr] (ui.east) -- node[label, above=-2pt] {WebSocket} (core.west);
  \draw[arr] (t.east) -- node[label, above=-2pt] {UDP} (net.west);

\end{tikzpicture}
\caption{What runs in and around a single peer. The Freenet Core is
a Rust process that hosts WebAssembly \emph{contracts} (whose state
is public and replicated across peers near the contract's ring
location) and \emph{delegates} (whose state is local to the device
and never leaves it). A \emph{user interface} is an HTML/JS/WASM
bundle delivered as contract state and run in the browser; it talks
to its local Core over a WebSocket. The Core's transport speaks UDP
with NAT hole-punching to neighbors on the rest of the network.}
\label{fig:peer-architecture}
\end{figure}


\subsection{Operations}

A small set of operations connect these artifacts. We describe them informally here; their detailed lifecycles appear in \cref{sec:routing}.

\begin{description}[leftmargin=1.4em, style=nextline, itemsep=2pt]
  \item[\textsf{PUT}] Publish a contract: insert or update its state at the peers near $\ell(\contract)$.
  \item[\textsf{GET}] Retrieve a contract's current state by key.
  \item[\textsf{UPDATE}] Submit an update to a contract. The update is merged into existing replicas using the contract's merge function and propagated outward through the subscription tree.
  \item[\textsf{SUBSCRIBE}] Register interest in future state changes to a contract.
  \item[\textsf{CONNECT}] Establish a peering relationship with the network as a new node.
\end{description}

\textsf{PUT}, \textsf{GET}, \textsf{UPDATE}, and \textsf{SUBSCRIBE} all route over the ring; their destination is $\ell(\contract)$. \textsf{CONNECT} is the bootstrap operation by which a fresh peer joins the topology.
